\documentclass[11pt]{article}
% Shared print layout. The web's provisional design and MathJax are independent.
\usepackage[utf8]{inputenc}
\usepackage[T1]{fontenc}
\usepackage[spanish,es-nodecimaldot]{babel}
\usepackage{lmodern}
\usepackage{amsmath,amssymb}
\usepackage{graphicx}
\usepackage{geometry}
\geometry{a4paper,left=20mm,right=20mm,top=25mm,bottom=22mm,headheight=16pt,headsep=8mm}
\usepackage{fancyhdr,lastpage}
\usepackage{enumitem}
\usepackage{xcolor}
\usepackage{pgfplots}
\pgfplotsset{compat=1.18}
\usepgfplotslibrary{fillbetween}
\usepackage[hidelinks,pdfencoding=auto]{hyperref}
\hypersetup{pdflang=es-ES,pdfauthor={},pdfsubject={Matemáticas Aplicadas a las Ciencias Sociales II}}
\pagestyle{fancy}
\fancyhf{}
\fancyhead[L]{\small Matemáticas Aplicadas a las Ciencias Sociales II}
\fancyhead[R]{\small 2.\textsuperscript{o} Bachillerato}
\fancyfoot[L]{\footnotesize Inferencia I | Edición revisada 1.0}
\fancyfoot[R]{\footnotesize Página \thepage\ de \pageref{LastPage}}
\renewcommand{\headrulewidth}{0.4pt}
\setlength{\parindent}{0pt}
\setlength{\parskip}{0.5em}
\setlist{itemsep=0.3em,topsep=0.3em}
\newcommand{\apartado}[1]{\subsection*{#1}}
\pgfplotsset{
  revision/.style={
    width=0.94\linewidth,height=7.5cm,
    axis lines=middle,grid=major,grid style={gray!18},
    tick label style={font=\small},label style={font=\small},
    xlabel={$x$},ylabel={$y$},samples=160,
    legend style={at={(0.5,-0.16)},anchor=north,draw=none,font=\small},
    clip=true
  }
}

\hypersetup{pdftitle={Inferencia I - Examen revisado}}
\begin{document}
\begin{minipage}{0.17\linewidth}\includegraphics[width=0.8\linewidth]{logo-al-andalus.png}\end{minipage}\hfill
\begin{minipage}{0.79\linewidth}{\Large\bfseries Inferencia I}\\[0.4em]
Duración: \textbf{55 minutos}. Puntuación: \textbf{10 puntos}.\end{minipage}

Apellidos y nombre: \hrulefill

Fecha: \hrulefill

Responde razonadamente y numera tus páginas. Se permite calculadora. Cada ejercicio incluye 0,5 puntos de expresión de procesos y métodos; no se añaden a su puntuación máxima. Se conservan los tiempos orientativos originales de 20, 20 y 10 minutos, dentro de los 55 minutos totales.


\section*{Ejercicio 1 \hfill 4 puntos} \textbf{(20')} La Delegación de Tráfico de una ciudad desea estudiar la influencia del uso del teléfono móvil en los accidentes de tráfico. Elegida una muestra aleatoria simple de 250 accidentes registrados el año pasado, se observó que 90 de ellos se produjeron por distracciones debidas al uso del móvil.

\begin{enumerate}[label=(\alph*)]

  \item (2,5 puntos)  Determine un intervalo de confianza al 97\% para estimar la proporción de accidentes de tráfico debidos al uso del móvil mientras se conduce.

  

  \item (1,5 puntos)  Usando la estimación anterior, calcule el tamaño mínimo que debe tener una muestra para estimar la proporción de accidentes con un error máximo del 5\% y un nivel de confianza del 99\%.

  

\end{enumerate}


\section*{Ejercicio 2 \hfill 4 puntos} \textbf{(20')} El peso de las ciruelas de una determinada variedad sigue una distribución Normal con media desconocida y desviación típica de 3 gramos. Se eligen al azar 25 ciruelas de esa variedad y se pesan, resultando un peso medio de 60 gramos.

\begin{enumerate}[label=(\alph*)]

  \item (2,5 puntos)  Calcule un intervalo al 95\% de confianza para estimar el peso medio de las ciruelas de esa variedad.

  

  \item (1,5 puntos)  Calcule el tamaño mínimo de la muestra que se ha de tomar, para que al estimar el peso medio de esa variedad de ciruelas con un nivel de confianza del 99\%, el error cometido sea inferior a 1 gramo.

  

\end{enumerate}

\clearpage
\section*{Ejercicio 3 \hfill 2 puntos}  \textbf{(10')} Una universidad desea realizar un estudio sobre los hábitos culturales del alumnado, dividiéndolos en cuatro estratos según la facultad a la que pertenecen: Ciencias Sociales, Ciencias de la Salud, Ingeniería y Humanidades, con tamaños de 4800, 2200, 3000 y 2000 estudiantes, respectivamente. Se realiza un muestreo estratificado con afijación proporcional y se han seleccionado 60 estudiantes de la Facultad de Ingeniería.

\begin{enumerate}[label=(\alph*)]
  \item (1 punto)  Determina el tamaño total de la muestra y cuántos estudiantes deben seleccionarse de cada estrato.
  
  

  \item (1 punto)  Explica la diferencia entre muestreo probabilístico y no probabilístico. Pon un ejemplo de cada tipo aplicado al estudio de hábitos culturales en la universidad, y reflexiona: \textquestiondown{}cuál crees que ofrece resultados más precisos? Justifica tu respuesta.

  
\end{enumerate}





\end{document}
